We cast simulator tuning as a bi-level optimization problem. Let $\mathcal{M}$ be the downstream perception model (e.g., MRISNet for ship segmentation \cite{ma2024mrisnet}) and $\mathcal{S}(\theta, e)$ be a physically grounded radar simulator parameterized by $\theta$ and conditioned on environment type $e \in \{\text{lake}, \text{river}, \text{coast}\}$. The objective is
\begin{equation}
    \theta^* = \arg\max_{\theta} \mathcal{R}\left(\mathcal{M}^*(\theta), \mathcal{D}_{real}\right)
\end{equation}
subject to inner-loop training
\begin{equation}
    \mathcal{M}^*(\theta) = \arg\min_{\omega} \mathcal{L}_{task}\left(\mathcal{M}_{\omega}, \mathcal{S}(\theta, e)\right).
\end{equation}

The outer-loop policy $\pi(\theta_{t+1} \mid s_t)$ proposes sequential updates to simulator parameters. State $s_t$ includes current $\theta$, environment label $e$, recent validation metrics, and trajectory history. Reward is
\begin{equation}
    r_t = \Delta \mathrm{mIoU}_{real} - \lambda_1\|\theta_t-\theta_{init}\|_2 - \lambda_2\,\mathbb{1}(\theta \notin \Omega_{phys}),
\end{equation}
where $\Omega_{phys}$ denotes physically plausible bounds.

The simulator models multipath and clutter using physically meaningful factors: speckle scale, attenuation, clutter floor, and multipath gain. To avoid unrealistic exploitation, each action is clipped to calibrated parameter ranges and penalized when violating physical plausibility. We use Soft Actor-Critic (SAC) for the outer loop because it supports stable learning in continuous action spaces with stochastic exploration.

The key design difference from prior adaptive randomization is that this policy is (i) sequential, (ii) environment-conditioned, and (iii) optimized on downstream real-task reward rather than simulator similarity alone.
